% ============================================================
%  Section 3: System and Optimizations
% ============================================================

\section{System}


\begin{frame}{Optimization Map}
  \centering
  \includegraphics[width=0.96\textwidth,height=0.78\textheight,keepaspectratio]{figures/tikz/optimization_legend.pdf}

  % \vspace{0.2em}
  % \small All three preserve the encoded probabilities. O1 and compressor-side O3 exploit that the input's next token is already known; O2 accelerates equivalent range-coder updates \parencite{Lam2015NumbaAL}.
\end{frame}

\begin{frame}{Experimental Setup}
  \small
  \begin{tabularx}{\textwidth}{lX}
    \toprule
    \textbf{Dimension} & \textbf{Configuration} \\
    \midrule
    Model & \texttt{Qwen/Qwen2.5-0.5B} \\
    Default hardware & Dual AMD EPYC Milan 7313, 512 GB DDR4, NVIDIA A100 40 GB SXM4 \\
    Datasets & \texttt{text8}, Python code from PyTorrent, Yelp Open Dataset records \\
    Default operating point & Batch size 256, context length 256, 100 retained tokens \\
    Metrics & Throughput in uncompressed MB/s; compression ratio as original/compressed size \\
    \bottomrule
  \end{tabularx}
\end{frame}

\begin{frame}{Three Datasets Cover Distinct Text Domains}
  \small
  \begin{tabularx}{\textwidth}{l>{\raggedright\arraybackslash}X>{\raggedright\arraybackslash}X}
    \toprule
    \textbf{Dataset} & \textbf{Content} & \textbf{Evaluation data} \\
    \midrule
    \texttt{text8} & English text extracted from Wikipedia & 100\,MiB \\
    PyTorrent & Python source code & First 100\,MiB \\
    Yelp Open Dataset & Structured business records & 150,346 records \\
    \bottomrule
  \end{tabularx}
  \vfill
  \tiny Datasets follow the evaluation corpus from prior work \parencite{Kipf2026WaitingTD}.
\end{frame}

\begin{frame}{Teacher-Forced Windows Amortize vLLM Scheduling}
  \centering
  \begin{tikzpicture}[scale=0.94,every node/.style={transform shape},
    req/.style={draw,rounded corners=2pt,minimum width=1.85cm,minimum height=0.7cm,align=center,font=\scriptsize},
    arr/.style={-{Stealth[length=5pt]},thick},
  ]
    \node[anchor=east,font=\small\bfseries] at (0.2,2.0) {one-token calls};
    \foreach \i/\x in {0/1.55,1/4.15,2/6.75,3/9.35} {
      \node[req,fill=red!9] (r\i) at (\x,2.0) {request + schedule\\score $t+\i$};
    }
    \draw[arr] (r0) -- (r1); \draw[arr] (r1) -- (r2); \draw[arr] (r2) -- (r3);

    \node[anchor=east,font=\small\bfseries] at (0.2,-0.25) {one window};
    \foreach \i/\x in {0/1.55,1/4.15,2/6.75,3/9.35} {
      \node[req,fill=blue!11] (s\i) at (\x,-0.25) {score $t+\i$\\force known token};
    }
    \draw[arr] (s0) -- node[above,font=\tiny] {KV state} (s1);
    \draw[arr] (s1) -- node[above,font=\tiny] {KV state} (s2);
    \draw[arr] (s2) -- node[above,font=\tiny] {KV state} (s3);
    \node[draw=mDarkTeal,very thick,rounded corners=2pt,fit=(s0)(s1)(s2)(s3),inner sep=6pt] {};
  \end{tikzpicture}

  % \vspace{0.7em}
  % \begin{block}{Why only vLLM?}
  %   vLLM's public path advances through a sampler and scheduler. Forcing lets one request traverse several known prefixes. Transformers already exposes a direct host-controlled token loop, so wrapping it in a generation request adds no benefit.
  % \end{block}
\end{frame}

\begin{frame}{Teacher-Forced Windows Amortize vLLM Scheduling}
  \centering
  \begin{tikzpicture}[
    box/.style={draw,rounded corners=2pt,minimum width=2.55cm,minimum height=0.9cm,align=center,font=\small},
    arr/.style={-{Stealth[length=5pt]},thick},
  ]
    \node[box,fill=blue!10] (score) at (0,1.3) {1. compute honest logits\\from current prefix};
    \node[box,fill=mLightBrown!15] (record) at (3.55,1.3) {2. record scores\\or target interval};
    \node[box,fill=green!12] (force) at (7.1,1.3) {3. force known $x_{t+1}$\\in the sampler};
    \node[box,fill=mDarkTeal!12] (advance) at (10.65,1.3) {4. append token\\advance KV cache};
    \draw[arr] (score) -- (record);
    \draw[arr] (record) -- (force);
    \draw[arr] (force) -- (advance);
    \draw[arr,rounded corners=5pt] (advance.south) -- ++(0,-0.75) -| node[below,pos=0.25,font=\scriptsize] {next internal decode step} (score.south);

    \node[draw,rounded corners=2pt,fill=black!4,align=center,font=\small,minimum width=5.35cm,minimum height=1.05cm] at (2.3,-1.15)
      {generation: sampled token\\chooses the next prefix};
    \node[draw,rounded corners=2pt,fill=mDarkTeal!8,align=center,font=\small,minimum width=5.35cm,minimum height=1.05cm] at (8.35,-1.15)
      {compression: observed file token\\chooses the next prefix};
  \end{tikzpicture}

  % \vspace{0.55em}
  % \small The score is captured \emph{before} forcing changes the sampler. KV caching avoids recomputing old keys and values, but each new position still runs the model.
\end{frame}

\begin{frame}{Windowing Raises vLLM Throughput $3.4\times$}
  \begin{columns}[T,onlytextwidth]
    \begin{column}{0.58\textwidth}
      \includegraphics[width=\textwidth]{\plotsdir/vllm_window_sweep_heatmap.pdf}
    \end{column}
    \hfill
    \begin{column}{0.37\textwidth}
      \vspace{0.55em}
      \begin{tikzpicture}
        \node[draw=mDarkTeal,very thick,rounded corners=3pt,fill=mDarkTeal!8,
          minimum width=4.2cm,minimum height=1.45cm,align=center]
          {\fontsize{28}{30}\selectfont\bfseries $3.43\times$\\[-0.05em]
           \small at the measured operating point};
      \end{tikzpicture}
      \vspace{0.7em}
      \small
      \begin{itemize}
        \item Gains concentrate at small and medium batches.
        \item Window 8 captures most of the benefit.
        \item Larger windows trade scheduling for memory.
      \end{itemize}
    \end{column}
  \end{columns}
  \vfill
  \tiny Qwen2.5-0.5B, vLLM, A100, \texttt{text8}; throughput in MB/s. Window 1 means one request per token.
\end{frame}

\begin{frame}[shrink=2]{Generation-Optimized Engines Expose the Wrong Interface}
  \centering
  \begin{tikzpicture}
    \begin{axis}[
      ybar,bar width=18pt,width=0.84\textwidth,height=0.45\textheight,
      ymin=0,ymax=0.038,ylabel={Compression throughput (MB/s)},
      symbolic x coords={Transformers,vLLM,TensorRT-LLM},xtick=data,
      tick label style={font=\small},label style={font=\small},
      grid=major,grid style={draw=black!13},scaled y ticks=false,
      yticklabel style={/pgf/number format/fixed,/pgf/number format/precision=3},
      nodes near coords,nodes near coords style={font=\scriptsize,/pgf/number format/fixed,/pgf/number format/precision=4},
    ]
      \addplot[fill=codecLLM!85,draw=codecLLM] coordinates {
        (Transformers,0.0344) (vLLM,0.0130) (TensorRT-LLM,0.0081)
      };
    \end{axis}
  \end{tikzpicture}

  \vspace{-0.2em}
  \small
  \begin{block}{High token serving throughput is not high compression throughput}
    Compression needs row-stable scores or intervals at every position. Scheduling, synchronization, and result reassembly become visible costs.
  \end{block}
  \vfill
  \tiny Qwen2.5-0.5B, batch 256, context 256, A100, \texttt{text8}; best comparable configuration per engine.
\end{frame}

% \begin{frame}{O1: Amortize the Request Boundary}
%   \centering
%   \begin{tikzpicture}[
%     reqbox/.style={draw,rounded corners=2pt,minimum width=1.8cm,minimum height=0.72cm,align=center,font=\scriptsize},
%     arr/.style={-{Stealth[length=5pt]},thick},
%   ]
%     \node[anchor=east,font=\small\bfseries] at (0.2,2.15) {baseline};
%     \foreach \i/\x in {0/1.5,1/4.2,2/6.9,3/9.6} {
%       \node[reqbox,fill=red!9] (r\i) at (\x,2.15) {request + schedule\\score $t+\i$};
%     }
%     \draw[arr] (r0) -- (r1);
%     \draw[arr] (r1) -- (r2);
%     \draw[arr] (r2) -- (r3);
%     \node[font=\scriptsize,text=red!70!black] at (5.55,1.38) {four request creations, schedules, and synchronizations};

%     \node[anchor=east,font=\small\bfseries] at (0.2,-0.35) {window $w=4$};
%     \node[reqbox,fill=blue!12] (s0) at (1.5,-0.35) {score $t$\\force $x_{t+1}$};
%     \node[reqbox,fill=blue!12] (s1) at (4.2,-0.35) {score $t+1$\\force $x_{t+2}$};
%     \node[reqbox,fill=blue!12] (s2) at (6.9,-0.35) {score $t+2$\\force $x_{t+3}$};
%     \node[reqbox,fill=blue!12] (s3) at (9.6,-0.35) {score $t+3$\\force $x_{t+4}$};
%     \draw[arr] (s0) -- node[above,font=\tiny] {KV state} (s1);
%     \draw[arr] (s1) -- node[above,font=\tiny] {KV state} (s2);
%     \draw[arr] (s2) -- node[above,font=\tiny] {KV state} (s3);
%     \node[draw=mDarkTeal,rounded corners=2pt,fit=(s0)(s1)(s2)(s3),inner sep=7pt,label={[font=\scriptsize,mDarkTeal]below:one vLLM request}] {};
%   \end{tikzpicture}

%   \vspace{0.8em}
%   \begin{block}{What O1 saves}
%     The same four model positions are evaluated. Windowing pays the fixed engine boundary once, rather than reducing the per-token transformer computation.
%   \end{block}
% \end{frame}


\begin{frame}{Multistreaming AC: Split the File into Row-Local Streams}
  \centering
  \begin{tikzpicture}[
    token/.style={draw,rounded corners=2pt,minimum width=0.75cm,minimum height=0.52cm,font=\scriptsize},
    box/.style={draw,rounded corners=2pt,minimum height=0.72cm,align=center,font=\small},
    arr/.style={-{Stealth[length=5pt]},thick},
  ]
    \node[font=\small\bfseries] at (1.45,3.2) {contiguous input rows};
    \node[font=\small\bfseries] at (6.0,3.2) {batched step $t$};
    \node[font=\small\bfseries] at (10.3,3.2) {independent coder states};

    \foreach \r/\y in {0/2.25,1/1.15,2/0.05} {
      \node[token,fill=mLightBrown!22] (seed\r) at (0,\y) {$x_{\r,1}$};
      \node[token,fill=blue!9] (a\r) at (0.95,\y) {$x_{\r,2}$};
      \node[token,fill=blue!9] (b\r) at (1.90,\y) {$\cdots$};
      \node[token,fill=blue!9] (c\r) at (2.85,\y) {$x_{\r,m_r}$};
      \node[box,fill=green!10,minimum width=2.15cm] (pred\r) at (6.0,\y) {$p_{\r,t}$ + true $x_{\r,t+1}$};
      \node[box,fill=mDarkTeal!12,minimum width=2.0cm] (coder\r) at (9.25,\y) {range coder $\r$};
      \node[box,fill=red!8,minimum width=1.6cm] (bits\r) at (11.85,\y) {bits $\r$};
      \draw[arr] (c\r) -- (pred\r);
      \draw[arr] (pred\r) -- node[above,font=\tiny] {$(\ell,h,T)$} (coder\r);
      \draw[arr] (coder\r) -- (bits\r);
    }

    \node[draw,rounded corners=2pt,fill=mLightBrown!10,align=center,font=\small,minimum width=4.1cm,minimum height=0.75cm] at (1.45,-1.25)
      {store the first token of each row\\to seed autoregressive decoding};
    \node[draw,rounded corners=2pt,fill=black!4,align=center,font=\small,minimum width=7.0cm,minimum height=0.75cm] at (9.1,-1.25)
      {parallel across rows; still sequential within each row};
  \end{tikzpicture}

  % \vspace{0.35em}
  % \small Row lengths and stream boundaries allow the independently encoded rows to be concatenated and reconstructed in original order.
\end{frame}

\begin{frame}{Multistream example: Two Rows, Two Bitstreams}
  \centering
  \begin{tikzpicture}[
    token/.style={draw,rounded corners=2pt,minimum width=0.62cm,minimum height=0.5cm,font=\small},
    box/.style={draw,rounded corners=2pt,minimum height=0.68cm,align=center,font=\small},
    arr/.style={-{Stealth[length=5pt]},thick},
  ]
    \node[anchor=east,font=\small\bfseries] at (-0.15,3.05) {input file};
    \foreach \x/\v in {0/A,0.72/B,1.44/A,2.16/C,2.88/A,3.60/B} {
      \node[token,fill=blue!9] at (\x,3.05) {\v};
    }
    \draw[very thick,mLightBrown] (1.80,2.65) -- (1.80,3.45);
    \node[font=\scriptsize,mLightBrown] at (1.80,2.43) {split into $B=2$ contiguous rows};

    \node[box,fill=mLightBrown!16,minimum width=1.65cm] (r0) at (0.45,1.25)
      {row 0\\\texttt{A B A}};
    \node[box,fill=mLightBrown!16,minimum width=1.65cm] (r1) at (0.45,-0.25)
      {row 1\\\texttt{C A B}};

    \node[box,fill=green!10,minimum width=3.55cm] (i0) at (4.1,1.25)
      {store seed \texttt{A}\\encode \texttt{B}: $[6,9)/10$\\encode \texttt{A}: $[0,7)/10$};
    \node[box,fill=green!10,minimum width=3.55cm] (i1) at (4.1,-0.25)
      {store seed \texttt{C}\\encode \texttt{A}: $[0,5)/10$\\encode \texttt{B}: $[2,9)/10$};

    \node[box,fill=mDarkTeal!13,minimum width=2.2cm] (c0) at (7.8,1.25)
      {coder state 0\\starts at $[0,255]$};
    \node[box,fill=mDarkTeal!13,minimum width=2.2cm] (c1) at (7.8,-0.25)
      {coder state 1\\starts at $[0,255]$};

    \node[box,fill=red!9,minimum width=1.65cm,font=\small\bfseries] (b0) at (10.75,1.25)
      {stream 0\\\texttt{11}};
    \node[box,fill=red!9,minimum width=1.65cm,font=\small\bfseries] (b1) at (10.75,-0.25)
      {stream 1\\\texttt{01}};

    \draw[arr] (r0) -- (i0);
    \draw[arr] (r1) -- (i1);
    \draw[arr] (i0) -- (c0);
    \draw[arr] (i1) -- (c1);
    \draw[arr] (c0) -- (b0);
    \draw[arr] (c1) -- (b1);

    \node[draw,rounded corners=2pt,fill=black!4,align=center,font=\scriptsize,
          text width=11.3cm,minimum height=0.82cm] at (5.45,-1.75)
      {payload = seeds \texttt{A,C} + stream lengths + stream 0 \texttt{11} + stream 1 \texttt{01}\\
       decode \texttt{A B A} and \texttt{C A B}, then join $\Longrightarrow$ \texttt{A B A C A B}};
  \end{tikzpicture}

  \vspace{0.15em}
  \scriptsize Toy alphabet $\{A,B,C\}$, frequency total $T=10$, and 8-bit coder state. Intervals come from each row's context-dependent model distribution; the displayed bits include the coder's terminating bit.
\end{frame}


\begin{frame}{Multistream makes the biggest difference}
  \centering
  \includegraphics[width=0.88\textwidth,height=0.66\textheight,keepaspectratio]{figures/tikz/bottleneck_shift.pdf}

  \vspace{0.15em}
  \begin{tikzpicture}
    \node[draw=codecBase,very thick,rounded corners=3pt,fill=red!6,
      minimum width=9.2cm,minimum height=0.8cm,align=center]
      {\textbf{97.1\%} of baseline wall time is arithmetic coding and score conversion};
  \end{tikzpicture}
  \vfill
  \tiny Qwen2.5-0.5B, Transformers, batch 256, A100, \texttt{text8}; end-to-end encode timing.
\end{frame}



\begin{frame}{Target-Interval Capture: Select Inside the Backend}
  \centering
  \begin{tikzpicture}[
    box/.style={draw,rounded corners=2pt,minimum height=0.78cm,align=center,font=\small},
    arr/.style={-{Stealth[length=5pt]},thick},
  ]
    \node[anchor=west,font=\small\bfseries] at (-1.6,2.75) {Baseline: dense scores cross the boundary};
    \node[box,fill=red!10,minimum width=3.2cm] (dense) at (0,1.75) {backend returns\\$B\times |V'|$ floats};
    \node[box,fill=red!10,minimum width=3.8cm] (host) at (4.4,1.75) {host softmax + integer counts\\prefix sum + target lookup};
    \node[box,fill=red!10,minimum width=2.6cm] (coder) at (8.8,1.75) {range-coder\\update};
    \draw[arr] (dense) -- (host);
    \draw[arr] (host) -- (coder);

    \node[anchor=west,font=\small\bfseries] at (-1.6,-0.1) {Target-Interval: select the known target before transfer};
    \node[box,fill=green!12,minimum width=3.2cm] (backend) at (0,-1.15) {backend softmax + counts\\prefix sum + target lookup};
    \node[box,fill=green!12,minimum width=3.8cm] (record) at (4.4,-1.15) {return $B$ records\\$(\ell,h,T,p_{target})$};
    \node[box,fill=mDarkTeal!12,minimum width=2.6cm] (coder2) at (8.8,-1.15) {range-coder\\update};
    \draw[arr] (backend) -- (record);
    \draw[arr] (record) -- (coder2);

    % \node[draw,rounded corners=2pt,fill=mLightBrown!12,align=left,font=\small,
    %       text width=10.9cm,inner sep=5pt] at (4.4,-2.75)
    %   {\textbf{Why integer counts?} Range coding updates an integer state, so encoder and
    %    decoder need the same cumulative bounds $(\ell,h,T)$. O3 moves this existing
    %    conversion; it does not add it.};
  \end{tikzpicture}
\end{frame}

\begin{frame}{Target-Interval Selection: One Row}
  \centering
  \begin{tikzpicture}[
    box/.style={draw,rounded corners=2pt,minimum height=0.72cm,align=center,font=\small},
    arr/.style={-{stealth[length=5pt]},thick},
  ]
    \node[font=\small\bfseries] at (1.75,3.0) {one model row};
    \node[box,fill=blue!10,minimum width=4.2cm] (probs) at (1.75,2.05)
      {alphabet $[a,b,c]$\\probabilities $[0.60,\,0.30,\,0.10]$\\known next file token: \texttt{b}};

    \node[font=\small\bfseries] at (7.25,3.0) {inside the inference backend};
    \node[box,fill=green!11,minimum width=3.75cm] (counts) at (6.2,2.05)
      {integer counts, $T=10$\\$[6,3,1]$\\CDF $[0,6,9,10]$};
    \node[box,fill=mLightBrown!16,minimum width=2.75cm] (select) at (10.25,2.05)
      {select \texttt{b}\\$[6,9)$};
    \draw[arr] (probs) -- (counts);
    \draw[arr] (counts) -- (select);

    \node[draw,rounded corners=2pt,fill=red!8,align=center,font=\small,
          minimum width=4.5cm,minimum height=1.05cm] (dense) at (2.25,-0.15)
      {\textbf{ordinary multistream handoff}\\transfer the dense row $[0.60,0.30,0.10]$\\host derives the target interval};
    \node[draw,rounded corners=2pt,fill=mDarkTeal!12,align=center,font=\small,
          minimum width=4.9cm,minimum height=1.05cm] (record) at (8.1,-0.15)
      {\textbf{target-interval handoff}\\transfer only $(\ell,h,T,p_{target})$\\$=(6,9,10,0.30)$};
    \draw[arr,red!65!black] (probs.south) -- (dense.north);
    \draw[arr,mDarkTeal] (select.south) -- (record.north);

    \node[box,fill=black!4,minimum width=6.6cm] (update) at (6.0,-1.75)
      {same row-local fastac coder: $[0,255]\rightarrow[153,229]$\\
       $low'=0+\lfloor 6\cdot256/10\rfloor$,\quad
       $high'=0+\lfloor 9\cdot256/10\rfloor-1$};
    \draw[arr] (dense.south) -- (update.north west);
    \draw[arr] (record.south) -- (update.north east);
  \end{tikzpicture}

  % \vspace{0.25em}
  % \small The integer range coder already needs $(\ell,h,T)$. O3 returns only that selected interval, removing the dense transfer and duplicate host work.
\end{frame}

\begin{frame}{Mutistream and Target Interval Raise Throughput $19\times$}
  \centering
  \begin{tikzpicture}
    \begin{axis}[
      xbar,bar width=11pt,width=0.83\textwidth,height=0.64\textheight,
      xmin=0,xmax=0.044,xlabel={Compression throughput (MB/s)},
      symbolic y coords={TensorRT target interval,TensorRT baseline,vLLM teacher-forced window,vLLM baseline,Transformers target interval,Transformers baseline},
      ytick=data,yticklabel style={font=\scriptsize},tick label style={font=\scriptsize},
      label style={font=\small},grid=major,grid style={draw=black!13},
      scaled x ticks=false,xticklabel style={/pgf/number format/fixed,/pgf/number format/precision=3},
      nodes near coords,nodes near coords style={font=\scriptsize,anchor=west,/pgf/number format/fixed,/pgf/number format/precision=4},
      enlarge y limits=0.10,
    ]
      \addplot[fill=codecLLM!85,draw=codecLLM] coordinates {
        (0.0018,Transformers baseline)
        (0.0347,Transformers target interval)
        (0.0031,vLLM baseline)
        (0.0116,vLLM teacher-forced window)
        (0.0020,TensorRT baseline)
        (0.0044,TensorRT target interval)
      };
    \end{axis}
  \end{tikzpicture}
  \vspace{-0.3em}

  \vfill
  \tiny Qwen2.5-0.5B, batch 256, context 128, A100, \texttt{text8}; relative gains use each engine's own baseline.
\end{frame}
